% Independent derivation of the complete coefficient-normalization extension.
% This file is an includable mathematical body; it performs no publication.
\section{Independent review: the complete ES coefficient extension}

The received derivation, \emph{Defining-prime continuation}, equations
(3)--(8), is the immediate source of the coefficient construction below.
The exact original receiving maps are those of the cumulative manuscript,
RD18--22 and GD14--17. All assertions made in this review are proved here.
In particular, a pullback to the normalization of the coefficient base is
not asserted to be the integral closure of the entire signed cover.

Put
\[
 R=\mathbb C[A,A^{-1},B,C],\qquad
 S=R[D]/(D^3-\tfrac{C}{5A}D^2+\tfrac{BC^2}{25A}D
                         -\tfrac{4C^4}{625A}),
 \quad
 N=R[P]/(P^3+A^{-1}P^2+(B/A)P+4C/(5A)).                 \tag{ETR1}
\]
Every occurrence of $A^{-1}$ has the literal domain $A\ne0$.
No restriction on $B$ or on $C$ is made until an inverse using $C^{-1}$
is explicitly introduced. Both displayed polynomials are monic; division
by them proves that the respective ordered bases over $R$ are
$(1,D,D^2)$ and $(1,P,P^2)$. Solving the relation for $C$ gives
\[
 C=-\tfrac54(AP^3+P^2+BP),\quad
 W=A^{-1}(B+P+AP^2),\quad
 B=AW-P-AP^2,
 \quad C=-\tfrac54APW,\quad D=\tfrac14AP^2W=-CP/5.       \tag{ETR2}
\]
Thus $N=\mathbb C[A,A^{-1},P,W]$. These formulas and their displayed
inverses prove the isomorphism, rather than only a parameter count.
The relation defining $S$ vanishes after substituting $D=-CP/5$.
The resulting inclusion has the complete matrix
\[
 K=\operatorname{diag}(1,-C/5,C^2/25),\qquad
 S\lhook\joinrel\longrightarrow N.                    \tag{ETR3}
\]
Since $R$ is a domain, this matrix is injective. It becomes invertible
over $R[C^{-1}]$, so the two domains have the same fraction field.
$N$ is finite over $S$: its three $R$-generators also generate it over
$S$, because $R\subset S$. A Laurent polynomial ring is integrally
closed: polynomial factorization over a field is unique, localization
preserves unique factorization, and an element $a/b$ with relatively
prime numerator and denominator satisfying a monic polynomial would
force every prime factor of $b$ to divide $a$. Hence $N$ is exactly the
normalization of $S$. Indeed, it is integral over $S$, and every element
of their common fraction field integral over $S$ is integral over $N$
by the same monic equation and consequently belongs to $N$.

The cokernel of ETR3, with its entire $R$-module structure, is
\[
 N/S=R/(C)\oplus R/(C^2).                                \tag{ETR4}
\]
Here and below the nonzero constants $5$ and $25$ are retained in the
matrices even though they are units for this cokernel calculation.
To calculate the algebra conductor, write an arbitrary element of $S$
uniquely as $a+CbP+C^2cP^2$, with $a,b,c\in R$. Write
$\alpha=A^{-1}$, $\beta=B/A$, $\gamma=4/(5A)$, so
$P^3=-\alpha P^2-\beta P-\gamma C$. The $P$ and $P^2$
coefficients of its product with $P$ are
$a-\beta C^2c$ and $Cb-\alpha C^2c$. Membership in $S$ forces
$a\in CR$ and $b\in CR$. The $P^2$ coefficient of its product with
$P^2$ is $a-\alpha Cb+(\alpha^2-\beta)C^2c$, which then forces
$a\in C^2R$. Conversely, $C^2N\subset S$ and is an ideal of $N$.
Since $1,P,P^2$ generate $N$, these tests prove the exact conductor
identities
\[
 (S:N)=C^2N=(C,D)^2\subset S.                            \tag{ETR5}
\]
The last equality follows from the three generators
$C^2,CD=-C^2P/5,D^2=C^2P^2/25$. It holds without assuming that
$B$ is invertible. At $C=D=0$ the distinguished parameter is retained
by $P(AP^2+P+B)=0$; assigning a value to $-5D/C$ would not define this
fiber.

\subsection{All twenty-four coordinates and the algebra conductor}
Keep the literal polynomial
\[
 h(r)=Ar^4+r^3+Br^2+Cr+D,\quad
 E=S[r,\sigma]/(h(r)/A,\ \sigma^2-h'(r)),\quad
 E_N=N\otimes_S E.                                     \tag{ETR6}
\]
Division first by the monic polynomial $h/A$ and then by the monic
polynomial in $\sigma$ proves that
$b=(1,r,r^2,r^3,\sigma,\sigma r,\sigma r^2,\sigma r^3)$ is an
$S$-basis of $E$ and an $N$-basis of $E_N$. Use the $R$-bases
$(b,Db,D^2b)$ and $(b,Pb,P^2b)$, retaining this order. Then
\[
 \iota=K\otimes I_8,
 \quad E_N/E=(R/(C))^8\oplus(R/(C^2))^8,
 \quad (E:E_N)=C^2E_N=(C,D)^2E.                         \tag{ETR7}
\]
For the conductor assertion, write $z=\sum_{i=1}^8z_i b_i\in E$.
The conditions $zP,zP^2\in E$ are exactly the eight separate conditions
$z_iP,z_iP^2\in S$, because $b$ is an $N$-basis of $E_N$. ETR5 makes
each $z_i$ belong to $C^2N$. Conversely, $E_N=E+EP+EP^2$, so those
conditions suffice for $zE_N\subset E$. This proves the conductor,
including its necessity, and does not invoke a claim that normalization
commutes with an arbitrary base change. The two algebras have rank
twenty-four over $R$, and the signed fiber still has all eight states
over each coefficient point. No pair of root labels or factor signs
has been identified.

\subsection{The distinguished action and its exact signed-root commutants}
In the indicated three-dimensional coefficient bases let
\[
 D_S=\begin{pmatrix}0&0&4C^4/(625A)\\1&0&-BC^2/(25A)\\0&1&C/(5A)\end{pmatrix},
 \quad
 P_N=\begin{pmatrix}0&0&-4C/(5A)\\1&0&-B/A\\0&1&-1/A\end{pmatrix}.
                                                               \tag{ETR8}
\]
On $C\ne0$ multiplication by $P=-5D/C$ on $E$ is
\[
 \mathcal M=M\otimes I_8,\quad
 M=\begin{pmatrix}0&0&-4C^3/(125A)\\-5/C&0&BC/(5A)\\0&-5/C&-1/A\end{pmatrix},
 \quad
 L=M^{-1}=\begin{pmatrix}-5B/(4C)&-C/5&0\\25/(4C^2)&0&-C/5\\-125A/(4C^3)&0&0\end{pmatrix}.
                                                               \tag{ETR9}
\]
Direct multiplication gives $ML=LM=I_3$ and
$\det M=-4C/(5A)$, $\det L=-5A/(4C)$. Also
$P_NK=KM$ and $D_S=-CM/5$. Thus $(P_N\otimes I_8)\iota
=\iota\mathcal M$ on the entire stated domain; the left-hand side
extends regularly to $C=0$.

Here are all operators needed to check that this distinguished action
retains the signed-root algebra. In zero-based indices put
\[
 J=\begin{pmatrix}0&0&0&0\\1&0&0&-C/A\\0&1&0&-B/A\\0&0&1&-1/A\end{pmatrix},
 \quad
 \mathcal R_S=I_3\otimes I_2\otimes J
       -A^{-1}D_S\otimes I_2\otimes E_{03},
 \quad
 \mathcal H_S=4A\mathcal R_S^3+3\mathcal R_S^2+2B\mathcal R_S+CI_{24},
 \quad
 \mathcal S_S=I_3\otimes E_{10}\otimes I_4
          +(I_3\otimes E_{01}\otimes I_4)\mathcal H_S .    \tag{ETR10}
\]
The first matrix is multiplication by $r$: its last column is exactly
the relation $r^4=-(r^3+Br^2+Cr+D)/A$. The second is multiplication
by $h'(r)$, and the last is multiplication by $\sigma$ in the ordered
even and odd parts. In particular,
\[
 A\mathcal R_S^4+\mathcal R_S^3+B\mathcal R_S^2+C\mathcal R_S
                      +D_S\otimes I_8=0,\quad
 \mathcal S_S^2=\mathcal H_S,\quad
 [\mathcal R_S,\mathcal S_S]=[\mathcal M,\mathcal R_S]
                         =[\mathcal M,\mathcal S_S]=0.   \tag{ETR11}
\]
The first identity follows column by column from the monic relation.
For the second, $\mathcal H_S$ is the same operator on the even and
odd copies, so multiplying the two by two matrix
$\left(\begin{smallmatrix}0&\mathcal H_S\\I&0\end{smallmatrix}\right)$
gives the claim in either copy. The commutators follow because every
coefficient in $\mathcal R_S$ is a polynomial in $D_S$, whereas
$M=-5D_S/C$; the even-odd matrices commute with both. Replace $D_S$
in ETR10 by $D_N=-CP_N/5$ to obtain $\mathcal R_N,\mathcal S_N$.
The equality $KD_S=D_NK$ proves, with no omitted correction terms,
\[
 \iota\mathcal R_S=\mathcal R_N\iota,\qquad
 \iota\mathcal S_S=\mathcal S_N\iota.                    \tag{ETR12}
\]
These equations are the exact relation with the root and signed-factor
operators; no assertion identifies these new operators with the
original conductor's multiplication by its spectral variable.

The pole residue is particularly explicit:
\[
 R_P=\left.(C\mathcal M)\right|_{C=0}
       =-5(E_{10}+E_{21})\otimes I_8,
 \quad R_P^2=25E_{20}\otimes I_8,
 \quad R_P^3=0,
 \quad (\operatorname{rank}R_P,\operatorname{rank}R_P^2,
             \operatorname{rank}R_P^3)=(16,8,0).         \tag{ETR13}
\]
These ranks follow from the independent images of the three coefficient
blocks: block zero maps onto block one, block one onto block two,
and block two to zero. There are eight Jordan blocks of size three.
This residue is the residue of the stated meromorphic action in the
coefficient lattice; the original conductor has not acquired this
nilpotent residue.

\subsection{Return to both fixed original metrics}
Let $\mathcal U,\mathcal W,T_A,\kappa$ have precisely their meanings
and original moments in RD18--22 and GD14. For a coefficient column
of $R(r)+\sigma S(r)$ define
\[
 F_U(R,S)=(T_A^{-1}R(S'-\kappa),T_A^{-1}S(S'-\kappa)),
 \quad F_W(R,S)=(R(S'-\kappa),S(S'-\kappa)),
 \quad G_8=T_A\oplus T_A,\quad F_W=G_8F_U.              \tag{ETR14}
\]
Their inverses first apply $T_A$ componentwise when necessary and
then substitute $S'=r+\kappa$. Thus the return is an actual linear
isomorphism on all eight coordinates. Let $\Gamma_U=F_U^*H_UF_U$
and $\Gamma_W=F_W^*H_WF_W$ denote the exact positive Gram matrices
from these original norms; neither is set equal to an identity.
On the full extension use
\[
 \mathcal F_U=I_3\otimes F_U,\quad
 \mathcal F_W=I_3\otimes F_W,\quad
 \Gamma_{24,U}=I_3\otimes\Gamma_U,\quad
 \Gamma_{24,W}=I_3\otimes\Gamma_W,\quad
 G_{24}=I_3\otimes G_8.                                 \tag{ETR15}
\]
The original center, Gamma masses and all inverse moments are retained
by these formulas; there are three copies of the exact eight-state
receiver. Define $\mathcal M_U=\mathcal F_U\mathcal M\mathcal F_U^{-1}$
and $\mathcal M_W=\mathcal F_W\mathcal M\mathcal F_W^{-1}$, with the same
definition for the root, sign and inclusion matrices. Substitution gives
\[
 G_{24}\mathcal M_U=\mathcal M_WG_{24},\quad
 G_{24}\mathcal R_U=\mathcal R_WG_{24},\quad
 G_{24}\mathcal S_U=\mathcal S_WG_{24}.                   \tag{ETR16}
\]
This gives the exact original conductor return and every commutant,
including the nilpotent residue. The norm on the coefficient block of
length three is the specified orthogonal direct sum, not a new metric
on any of the original eight-state spaces.

\subsection{Finite singular values and every exterior inverse norm}
For $A,C\ne0$ and arbitrary complex $B$, put
\[
 u=\begin{pmatrix}-5B/(4C)\\25/(4C^2)\\-125A/(4C^3)\end{pmatrix},\quad
 t=|C|^2/25,\quad r_u=u^*u,\quad q=|u_3|^2
      ={15625|A|^2\over16|C|^6}.
                                                               \tag{ETR17}
\]
Here $u_3$ denotes the third entry of the displayed column. Since
$L=(u,-Ce_1/5,-Ce_2/5)$, one has
$LL^*=uu^*+t\operatorname{diag}(1,1,0)$. The nonzero vector in
$\operatorname{span}(e_1,e_2)$ orthogonal to $(u_1,u_2,0)$ is an
eigenvector of eigenvalue $t$. On its orthogonal complement the
trace is $r_u+t$ and the determinant is $tq$: in the orthonormal basis
consisting of $(u_1,u_2,0)/\sqrt{|u_1|^2+|u_2|^2}$ and $e_3$, this
follows by multiplying the full two by two Hermitian matrix. Thus
\[
 \lambda_\pm={r_u+t\pm\sqrt{(r_u+t)^2-4tq}\over2},\quad
 \lambda_-={2tq\over r_u+t+\sqrt{(r_u+t)^2-4tq}},\quad
 (\alpha,\beta,\chi)=(\sqrt{\lambda_+},\sqrt t,\sqrt{\lambda_-}).
                                                               \tag{ETR18}
\]
The three symbols on the right are singular values, not ES root labels.
The characteristic polynomial evaluated at $t$ is
$-t(|u_1|^2+|u_2|^2)<0$, because $u_2\ne0$. Its positive roots therefore
satisfy $\lambda_+>t>\lambda_->0$, proving their order and the strict
positivity of the displayed square root. These are exact finite formulas.

For any positive original Gram $\Gamma$, conjugation by
$I_3\otimes\Gamma^{1/2}$ sends $L\otimes I_8$ to itself. This is an
isometry from the literal Gram $I_3\otimes\Gamma$ to the Euclidean
calculation and proves that its twenty-four singular values in either
original receiver are
\[
 \underbrace{\alpha,\ldots,\alpha}_{8},\quad
 \underbrace{\beta,\ldots,\beta}_{8},\quad
 \underbrace{\chi,\ldots,\chi}_{8}.                     \tag{ETR19}
\]
The cancellation here is a proved tensor identity, not an assumption
that the original Gram is the identity. In particular only the middle
eight values are exactly $|C|/5$. The last eight are strictly smaller
for every $A,C\ne0$.

For $0\le j\le24$, with an exterior power of degree zero assigned
norm one, define $a_j=\min(j,8)$, $b_j=\min(\max(j-8,0),8)$ and
$c_j=\max(j-16,0)$. The complete exterior inverse answer is
\[
 \|\mathop{\bigwedge}\nolimits^j\mathcal M^{-1}\|
       =\alpha^{a_j}\beta^{b_j}\chi^{c_j}.               \tag{ETR20}
\]
To prove this, apply an orthonormal singular basis in both the domain
and codomain; its exterior basis is again orthonormal, and the exterior
matrix is diagonal with every product of $j$ distinct singular values.
Its largest entry is exactly the ordered product in ETR20.

There is also a finite explicit two-sided bound with the exact leading
constants. Set
\[
 k={125|A|\over4},\quad
 \delta={|C|^2\over25|A|^2}
       +{|B|^2|C|^4\over625|A|^2}
       +{16|C|^8\over390625|A|^2},\quad
 Q_j=(k|C|^{-3})^{\min(j,8)}(|C|/5)^{\max(j-8,0)},\quad
 e_j=\min(j,8,24-j).
                                                               \tag{ETR21}
\]
Because the largest eigenvalue of $uu^*+t\operatorname{diag}(1,1,0)$
is at least $r_u\ge q$ and at most $r_u+t=q(1+\delta)$, and
$\lambda_+\lambda_-=tq$, one has, on the entire finite domain,
\[
 Q_j\ \le\ \|\mathop{\bigwedge}\nolimits^j\mathcal M^{-1}\|
       \ \le\ Q_j(1+\delta)^{e_j/2}.                    \tag{ETR22}
\]
Indeed, writing $\theta=\alpha/\sqrt q\in[1,\sqrt{1+\delta}]$ gives
$\chi=\beta/\theta$ and the exact identity
$\|\bigwedge^j\mathcal M^{-1}\|=Q_j\theta^{a_j-c_j}$;
$a_j-c_j=e_j$ in all three ranges. This proves ETR22 without discarding
any direction. At fixed $A\ne0$ and fixed $B$, $\delta\to0$ as $C\to0$,
so the complete leading terms are
\[
 \|\mathop{\bigwedge}\nolimits^j\mathcal M^{-1}\|
 \sim
 \begin{cases}
 k^j|C|^{-3j},&0\le j\le8,\\
 k^8 5^{-(j-8)}|C|^{j-32},&8\le j\le24.
 \end{cases}\qquad
 \|\mathop{\bigwedge}\nolimits^{24}\mathcal M^{-1}\|
       =\left|{5A\over4C}\right|^8.                     \tag{ETR23}
\]
If $A$ varies, the limit requires $A$ to have a nonzero limiting value;
the exact ETR18--22 remain valid without this limiting hypothesis.

Finally retain the original conductor in a cross-receiver calculation.
Let $g_1,\ldots,g_8$ be the singular values of the literal
$G_8=T_A\oplus T_A$ with its given original source and target norms.
Each of the four singular values of $T_A$ occurs twice. In coefficient
coordinates these are the singular values of
$\Gamma_W^{1/2}\Gamma_U^{-1/2}$, because both coefficient columns
describe the same transported vector under ETR14. The full operator
$G_{24}\mathcal M_U^{-1}$, from the source receiver to the target
receiver, consequently has exactly the multiset
\[
 \{\alpha g_i,\ \beta g_i,\ \chi g_i:1\le i\le8\}.    \tag{ETR24}
\]
An orthonormal singular basis for each tensor factor proves this by
direct multiplication. Every cross-receiver exterior norm is the product
of the largest $j$ members of this displayed, entirely retained
multiset. Thus ETR24 retains the original moments and Gamma factors in
the cross-receiver estimate. It does not replace them by one. The fixed
conductor $G_{24}$ itself remains invertible; the pole and the conductor
ideal in ETR7 and ETR13 belong to the explicitly introduced coefficient
extension and action. No order jump of the original $T_A$, and no
statement settling RH, follows from these calculations alone.
\subsection{The inverse of the full source-to-target comparison}
For the actual source-to-target operator, retain its order:
\[
 \mathcal B_C=\mathcal F_W(M\otimes I_8)\mathcal F_U^{-1}
              =G_{24}\mathcal M_U,\qquad
 \mathcal B_C^{-1}=\mathcal F_U(L\otimes I_8)\mathcal F_W^{-1}
              =\mathcal M_U^{-1}G_{24}^{-1}.             \tag{ETR25}
\]
Both products are the identity by ETR9 and the exact receiver inverses.
Let $\gamma_1\ge\cdots\ge\gamma_8>0$ be the singular values of the
original $G_8^{-1}$, from its target metric to its source metric. They
are the original four singular values of $T_A^{-1}$ each repeated
twice. Applying the definition of the metric adjoint to ETR25 gives
\[
 (L^*L)\otimes(\Gamma_W^{-1}\Gamma_U),\qquad
 \operatorname{Sing}(\mathcal B_C^{-1})
     =\{\alpha\gamma_j,\beta\gamma_j,\chi\gamma_j:1\le j\le8\}.
                                                               \tag{ETR26}
\]
The second factor is self-adjoint and positive for the $\Gamma_W$
inner product. It is similar to the positive Hermitian matrix
$\Gamma_W^{-1/2}\Gamma_U\Gamma_W^{-1/2}$ and hence has a full
orthonormal singular basis for that metric. Its eigenvalues are exactly
$\gamma_j^2$, as follows by pulling back the unchanged map $G_8^{-1}$
through ETR14. Tensoring the orthonormal eigenbases proves ETR26,
including all multiplicities. Every finite exterior inverse norm is
therefore the product of the largest $j$ members of this explicit list.

Define the decreasing sixteen-element list
$\widehat\gamma_{2i-1}=\widehat\gamma_{2i}=\gamma_i$.
At fixed $A\ne0$, fixed $B$ and fixed original conductor, the smallest
of the eight large values divided by the largest of the sixteen small
ones is at least $\alpha\gamma_8/(\beta\gamma_1)$, which tends to
infinity. The large and small groups thus separate for all sufficiently
small $|C|>0$. Within the large group their order is exactly that of
the $\gamma_i$. Dividing the small group by $|C|/5$ gives a list
converging entry by entry to the duplicated list $\widehat\gamma$,
because $\chi/\beta\to1$. Sorting is continuous for finite real lists:
if every entry changes by at most $\varepsilon$, every order statistic
changes by at most $\varepsilon$, since the counts above any threshold
can change only in its $\varepsilon$ neighborhood. This argument also
covers tied values. Consequently, with $k=125|A|/4$,
\[
 \|\mathop{\bigwedge}\nolimits^j\mathcal B_C^{-1}\|
 \sim\begin{cases}
 k^j(\prod_{i=1}^{j}\gamma_i)|C|^{-3j},&0\le j\le8,\\
 k^8 5^{-(j-8)}(\prod_{i=1}^{8}\gamma_i)
      (\prod_{i=1}^{j-8}\widehat\gamma_i)|C|^{j-32},&8\le j\le24.
 \end{cases}                                                \tag{ETR27}
\]
Empty products are one. In degree twenty-four, determinant
multiplicativity gives the exact finite equality
\[
 \|\mathop{\bigwedge}\nolimits^{24}\mathcal B_C^{-1}\|
   =\left|\frac{5A}{4C}\right|^8 |\det A_{3,s}|^{-6},\qquad
 |\det A_{3,s}|=
 \left|\frac{(-i)^v\mu_v}{v!}\right|^4
       \prod_{i=0}^{3}\frac{\rho_{v+i}^{(s)}}{\rho_i^{(s)}}.
                                                               \tag{ETR28}
\]
The last formula is the full original FC23 determinant, which is
nonzero and has all its original weights. The exponent six is forced
by the six original four-dimensional conductor components, equivalently
by $(\prod_i\gamma_i)^3=|\det A_{3,s}|^{-6}$. Thus the total exterior
estimate preserves the entire original conductor and every one of the
twenty-four extension directions.

\subsection{Independent check of the coefficient differential and its fibers}
For completeness the full differential of ETR2, with target order
$(A,B,C,D)$ and source order $(A,P,W)$, is
\[
 J_\nu=\begin{pmatrix}
 1&0&0\\ W-P^2&-1-2AP&A\\-5PW/4&-5AW/4&-5AP/4\\
 P^2W/4&APW/2&AP^2/4
 \end{pmatrix}.
                                                               \tag{ETR29}
\]
Its maximal minors in row orders $ABC,ABD,ACD,BCD$ are respectively
$5A(2AP^2+AW+P)/4$, $-AP(2AP^2+2AW+P)/4$,
$5A^2P^2W/16$, and $-5A^2P^4W/16$, by direct expansion.
Their common zero locus is the union $L_0=\{P=W=0\}$ and
$L_1=\{W=0,P=-1/(2A)\}$: the third minor first forces $P=0$ or
$W=0$, and the first two then give exactly these alternatives. The
first row and the nonzero $A$ entry in the second row show rank two,
and direct substitution gives kernel vectors $(0,A,1)$ and $(0,1,0)$.
The full signed fibers are
\[
 L_0:\quad h=r^3(Ar+1),\quad
 E_0=\mathbb C[r,\sigma]/(r^3,\sigma^2-3r^2)
       \ \times\ \mathbb C_{\sigma=i/A}\ \times\ \mathbb C_{\sigma=-i/A};
 \qquad
 L_1:\quad h=Ar^2(r+1/(2A))^2,\quad
 E_1=\prod_{r_0=0,-1/(2A)}
 \mathbb C[\epsilon,\sigma]/(\epsilon^2,\sigma^2-\epsilon/(2A)),
 \quad\epsilon=r-r_0.                                      \tag{ETR30}
\]
For $L_0$, reducing $h'$ modulo $r^3$ gives $3r^2$; evaluating at
$-1/A$ gives $-1/A^2$, with exactly the two stated signs. For $L_1$,
expansion at either original root and reduction modulo $\epsilon^2$
gives $\epsilon/(2A)$. The factors are pairwise coprime because
$A\ne0$, so the ordinary Bezout identity gives the indicated direct
products. Their dimensions are respectively $6+1+1$ and $4+4$.
No signed state is deleted by either differential rank loss.

\subsection{A uniform finite bound for the full cross-receiver inverse}
Keep the complete original positive list $\gamma_1,\ldots,\gamma_8$
from ETR26. Write $a=125|A|/(4|C|^3)$, $b=|C|/5$, and use the
unchanged $\delta$ of ETR21. Let $L_j$ be the product of the largest
$j$ entries of the twenty-four-element multiset
$\{a\gamma_i,b\gamma_i,b\gamma_i:1\le i\le8\}$, with $L_0=1$.
Then the following bounds hold for every $A,C\ne0$, every complex
$B$, and every exterior degree, without assuming any ordering of
$a$ and $b$:
\[
 L_j\ \le\ \|\mathop{\bigwedge}\nolimits^j\mathcal B_C^{-1}\|
       \ \le\ L_j(1+\delta)^{\min(j,8,24-j)/2},
       \qquad 0\le j\le24.                              \tag{ETR31}
\]
For a full proof put $\theta=\alpha/a$. ETR18--22 give
$1\le\theta\le\sqrt{1+\delta}$ and the exact actual triple
$(\alpha,\beta,\chi)=(a\theta,b,b/\theta)$.
The finite order proved in ETR18 gives $\alpha>b$; the Rayleigh
quotient used in ETR21 gives $\alpha\ge a$. Thus $\alpha$ is at
least both $a$ and $b$. In one fixed $i$ block, the largest one-entry
actual product is $\alpha\gamma_i$, at least
$\max(a,b)\gamma_i$, which is the largest baseline one-entry
product. The largest two-entry actual product is $\alpha b\gamma_i^2$;
its difference from each candidate baseline two-entry product is
$b(\alpha-a)\gamma_i^2$ or $b(\alpha-b)\gamma_i^2$, both nonnegative.
For zero entries both products are one; for three entries both are
$ab^2\gamma_i^3$, since the $\theta$ and $\theta^{-1}$ factors cancel.
Choose a subset of baseline entries attaining $L_j$ and retain its
number of entries in each of the eight blocks. Replacing those entries
in each block by that many largest actual entries can only increase
the product, by the four cases just proved. The resulting actual
$j$-element product is bounded above by its largest such product,
$\|\bigwedge^j\mathcal B_C^{-1}\|$. This proves the lower bound.

For the upper bound, assign to each actual entry the identically
indexed baseline entry. Their ratios are respectively $\theta$, $1$
and $\theta^{-1}$ in every block. For a selected $j$-element subset
let $n_+,n_0,n_-$ count these three types. Its exponent is
$e=n_+-n_-$. The inequalities $e\le j$ and $e\le8$ follow from
$n_++n_0+n_-=j$ and $n_+\le8$. The full list has total exponent zero;
its complement contains $24-j$ entries, each with exponent at least
$-1$, so $e$ is at most $24-j$. Hence
$e\le\min(j,8,24-j)$ for every subset. Since $\theta\ge1$, its
actual product is at most its baseline product times
$\theta^{\min(j,8,24-j)}$, and the baseline product is at most
$L_j$. Maximizing over the subsets proves the upper bound using
$\theta\le\sqrt{1+\delta}$. This proof retains coincident conductor
singular values and any changes in their order. At $j=24$ the exponent
is zero and both bounds equal the exact determinant ETR28.
